%% 第五章 结论与展望

\chapter{结论与展望}
\chaptermark{结论与展望}

\section{本文工作总结}

本文以无人机在非定常大气环境中的节能飞行为应用背景，
将滑翔问题建模为自推进体在非定常流场中的点对点迁移问题，
结合活性物质物理理论与深度强化学习方法，
系统研究了基于Soft Actor-Critic算法的能耗优化迁移策略。
主要完成的工作和结论如下：

(1) 构建了流场与自推进体一体化的数值研究框架。
以二维非定常双涡流模型为理论验证环境，建立了自推进体的运动学方程，
定义了推进功率与能量消耗的量化方法，为后续的策略学习与性能评估
提供了完整的数学模型。
该框架的突出特点在于其模块化设计——流场模型、运动学模型、
能耗定义与评价指标体系之间相互解耦，
使得后续研究者可以便捷地替换流场模型（如从双涡流切换至Rayleigh-B\'{e}nard对流）
或运动学模型（如引入惯性项）而不必重构整体的算法与评估流程。

(2) 设计并实现了面向能耗优化迁移任务的强化学习算法。
基于最大熵框架的SAC算法，设计了状态奖励、能量奖励与
时间奖励三元复合奖励函数，以最小能耗和成功到达为核心
优化目标；配合直线基准策略实现公平的节能效果评估。
三元复合奖励函数的设计是本文在算法层面的核心贡献之一：
通过将空间约束、能量约束和时间约束编码为标量奖励信号，
使得强化学习智能体无需复杂的多目标优化框架
即可自主平衡迁移效率与能耗之间的竞争关系。

(3) 通过数值仿真验证了强化学习策略的节能有效性。
经过训练的智能体能够自主习得“借流而行”的迁移策略——
先顺应涡旋流线运动以获取流场对流的“助力”，
再调整方向切入目标区域。其能量消耗仅为直线基准策略的约五分之一，
节能效果十分显著。

(4) 分析了涡旋振幅对迁移策略形态的影响规律。
当 $A$ 过小时，智能体无法有效利用流场对流，
迁移效率低下；但当流场强度过大时，策略同样失效——
流场对流远超自主推进能力，智能体被吹离有效区域。
智能体迁移策略对横向振幅和角频率的变化均具有较强的鲁棒性。

(5) 验证了训练策略的收敛性与泛化能力。
训练奖励在约100000步后趋于收敛，策略从随机探索逐步演化为
稳定的最优行为。随机起终点测试表明，策略能够将习得的流场结构知识
泛化至未见过的起终点组合，验证了学习结果的可迁移性。

在方法层面，可引入偏好驱动的多目标强化学习框架，
通过调整奖励函数中能耗项与时间项的权重参数，
生成不同偏好设置下的Pareto前沿策略集，
供工程决策者根据具体任务需求
灵活选用极度节能或快速到达等不同策略。

这一泛化能力根植于SAC算法所学到的并非特定起终点对之间的固定路径映射，
而是双涡流流场的通用几何特征——
如涡旋中心的相对位置、异宿轨道的走向以及流线曲率的空间分布。
换言之，智能体学习到的是流场的结构知识而非导航记忆，
前者具有空间平移不变性，
因此可以在未经训练的起终点组合上直接复用。

从实际应用角度看，这一泛化能力意味着：
无人机在实际部署前，可以在有限的仿真场景中进行训练，
随后无需针对每一组新的起终点坐标重新训练即可直接应用，
大幅降低了在线计算成本和部署门槛。

(6) 通过余弦相似度分析揭示了“借流而行”的物理机制。
智能体推进方向与流场方向保持约60度的平均夹角，
表明智能体在“顺流省力”与“横向寻路”之间实现了最优平衡——
这一发现将节能行为从经验观察提升至物理机理层面。

\section{创新点}



(1) 研究范式创新：将无人机滑翔节能问题纳入活性物质物理与强化学习交叉的研究范式，
建立了从宏观飞行器到自推进颗粒的跨尺度建模框架。
不同于传统飞行动力学仅从气动优化的角度处理滑翔问题，
本文从非平衡统计物理的视角出发，
将无人机视为非平衡能量场中的自驱动活性颗粒，
揭示了滑翔行为与活性物质输运之间的深层理论关联。
这一范式转换使得活性物质理论中关于能量转化效率、
非平衡稳态和最优输运策略的丰富成果
可以被系统性地引入无人机滑翔研究领域。

(2) 方法学创新：在能耗评估中引入直线基准策略的时间对齐控制变量设计，
确保了节能对比的公平性；引入余弦相似度指标，
从矢量对齐角度量化了借流而行的物理机制。
与现有研究中仅比较最终能耗值的方法不同，
本文的时间对齐设计排除了迁移耗时的干扰因素，
使能耗差异纯粹归因于路径与策略选择；
余弦相似度指标则为节能行为提供了连续时间分辨率的微观解释，
将结果分析从最终数值对比推进到了
每一时间步的决策机理层面。
此外，三元复合奖励函数的设计——
将空间约束、能量约束和时间约束统一在标量奖励信号中——
为多目标强化学习在节能迁移问题中的应用提供了可复用的设计范式。

\section{研究展望}

本文的工作仍存在若干可进一步拓展的方向：

(1) 流场复杂度的提升。本文以二维双涡流为理论验证环境，
流场虽然具有非定常和涡旋结构特征，但本质上仍是低维且周期性的。
将研究拓展至瑞利-伯纳德热对流的湍流环境，
可以更真实地模拟大气对流层的流场条件，
对验证强化学习方法在真实湍流中的适用性具有重要意义。
具体而言，可将本文的双涡流模型替换为Rayleigh-B'{e}nard对流的DNS数值模拟环境，
在Rayleigh数$Ra = 10^6$至$10^9$的范围内研究智能策略
在不同湍流强度下的演化规律。
湍流热对流中天然存在的大尺度环流与热羽流结构
为无人机利用大气热上升气流提供了更逼真的物理类比。

(2) 颗粒惯性的引入。本文假设自推进体为无惯性颗粒，
即速度响应瞬时完成。然而，真实飞行器具有一定的惯性，
其运动响应存在时间延迟，且惯性与流场涡结构的相互作用
可能导致新的迁移策略和能耗特征。
在实现路径上，可在运动学方程中引入Maxey-Riley方程
或简化的Stokes拖曳模型来描述有限惯性颗粒的动力学响应。
通过改变Stokes数$St = \tau_p / \tau_f$（其中$\tau_p$为颗粒弛豫时间，
$\tau_f$为流场特征时间），系统研究惯性效应对智能体
捕获利流结构能力和能耗特征的影响。

(3) 多智能体协同迁移。在无人机集群任务中，
多架飞行器之间的协同路径规划与能耗分配是一个复杂的最优控制问题。
将本文的单智能体框架推广至多智能体强化学习，
可探索群体“借流”的协同节能策略。
具体技术路线可基于Multi-Agent SAC (MASAC)或QMIX等
集中训练-分布执行的多智能体强化学习框架，
研究无人机集群在共享流场中的编队维持、
碰撞避免与协同借流问题。
特别值得关注的是，多智能体系统中可能出现
领航-跟随的流场利用策略——
前排无人机为后排无人机创造有利的尾流条件，
从而实现超出单机滑翔理论极限的集群整体能效。

(4) 真实飞行器动力学约束。自推进体点状模型忽略了飞行器的
姿态动力学、气动约束和最小转弯半径等实际限制。
未来可将本文的简化模型与更真实的飞行动力学模型结合，
增强策略在实际飞行任务中的可部署性。
可行的方法包括：首先在六自由度飞行动力学模型中
引入最小转弯半径、最大滚转角和最大过载等约束，
然后通过约束强化学习或安全强化学习方法
将本文习得的点状颗粒最优轨迹
映射为满足动力学约束的可飞行轨迹。
此外，可探索将本文策略作为高层路径规划器、
底层控制器负责轨迹跟踪的分层控制架构。

(5) 从仿真到实验的验证。本文的结论基于数值仿真，
下一步可考虑在水洞或风洞实验中构造可控自推进体模型，
在物理环境中验证强化学习策略的实际节能效果。
在风洞实验方面，可设计携带惯性测量单元和微型推进器的
中性浮力球体模型，在可控振荡的二维水洞流场中
部署本文训练的强化学习策略，
通过粒子图像测速（PIV）系统同步记录
智能体的运动轨迹和流场结构，
定量对比仿真结果与实验数据的吻合程度。


综上所述，本文从活性物质物理与人工智能的交叉视角出发，
以双涡流流场为理论验证平台，
系统论证了强化学习方法在自推进体能耗优化迁移问题中的有效性。
研究发现的核心启示在于：
在非平衡能量环境中，
最优迁移策略并非简单地加速或对抗，
而是通过感知并顺应环境能量梯度来实现高效输运。
这一结论不仅适用于无人机滑翔场景，
也为更广泛的非平衡输运问题——
从微型医疗机器人在血管流中的导航
到海洋滑翔机在洋流中的长航时部署——
提供了统一的理论视角和方法论框架。
随着活性物质物理理论与人工智能方法的持续发展，
二者交叉融合所释放的学术潜力与应用价值值得期待。